\chapter*{General Introduction}
\label{ch:intoGen}

Reinforcement Learning (RL) is a paradigm of machine learning in which an 
agent learns optimal behaviour by interacting with an environment, receiving 
numerical reward signals, and adjusting its decision-making policy over time. 
Unlike supervised learning, no labelled data is required: the agent discovers 
useful strategies purely through trial and error. This property makes RL 
particularly well-suited to problems where the optimal solution is unknown in 
advance, such as playing games, robotic control, and autonomous navigation.

This project investigates how deep reinforcement learning algorithms can be 
applied to a physically realistic, multi-agent football environment. The 
environment simulates four embodied agents — two per team — competing in 
a 2-versus-2 soccer match inside the MuJoCo physics engine. The agents must 
learn, without any human guidance, to chase the ball, cooperate with their 
teammate, and score goals against their opponents. This problem is challenging 
because the state space is continuous and high-dimensional, rewards are sparse, 
and the presence of multiple simultaneously learning agents makes the 
environment non-stationary.

The report traces the full learning journey: from the theoretical foundations 
of Q-Learning and the Bellman equation, through the advances introduced by Deep 
Q-Networks, and finally to the state-of-the-art Proximal Policy Optimisation 
algorithm used to train the multi-agent soccer system. Both implementation and 
evaluation are covered in detail, with results demonstrated through learning 
curves and statistical analysis.

\section*{Organization of the report}
\label{sec:intro_org}

This report is organised into seven chapters. 
Chapter~\ref{ch:into} introduces the host organisation, the problem statement, 
and the project objectives. 
Chapter~\ref{ch:lit_rev} reviews the theoretical and algorithmic foundations 
of reinforcement learning covered in the course. 
Chapter~\ref{ch:method} describes the methodology, including the environment 
design, reward shaping strategy, and training configuration. 
Chapter~\ref{ch:results} presents the quantitative results obtained at each 
stage of the pipeline. 
Chapter~\ref{ch:evaluation} provides a critical discussion and analysis of 
the findings. 
Chapter~\ref{ch:conclusion} draws conclusions and outlines directions for 
future work. 
Chapter~\ref{ch:reflection} reflects on the personal learning experience 
gained through this project.
